%%%%%
%
\section{Correction de vitesse sur l'humidit� relative}

\itemstyle{Op\'erateur~:} \operfmt{Calc\_Static\_RH\_From\_Dynamic\_RH\_V1}
\label{OP-Calc_Static_RH_From_Dynamic_RH_V1}
\index{Calc\_Static\_RH\_From\_Dynamic\_RH\_V1}

%%%%%OBJET%%%%
\itemstyle{Objet~:}
Correction des effets de compression et de
r�chauffement qui affectent la mesure d'humidit� relative
sur plate-forme mobiles.

%%%%%ARGUMENTS%%%%
\bigskip
%%%%%ENTREES%%%%
\begin{tabarg}
\multicolumn{3}{l}{\titretab {Entr\'ees}}\\
\hline
$HU_{dyn}$ & Vecteur, \texttt{float} & Mesure d'humidit� relative dans des
conditions de vitesse/surpression (\%)\\
$T_s$ & Vecteur, \texttt{float} & Temp�rature statique de l'air (�C ou K)\\
$T_t$ & Vecteur, \texttt{float} & Temp�rature totale (�C ou K)\\
\end{tabarg}

\medskip
%%%%%SORTIES%%%%
\begin{tabarg}
\multicolumn{3}{l}{\titretab {Sorties}}\\
\hline
$HU_{stat}$ & Vecteur, \texttt{float} & Humidit� relative corrig�e des effets de 
la vitesse (\%)\\
\end{tabarg}


%%%%%FORMULE%%%%
\itemstyle{Formule}

%% Calc_Static_RH_From_Dynamic_RH_V1
Si besoin est, $T_s$ et $T_t$ sont convertis en kelvins, puis
\[
HU_{stat} = 
\frac{e_w(T_t)}{e_w(T_s)}
\left(\frac{T_s}{T_t}\right)^{\frac{\gamma}{\gamma-1}}
% \gamma = 1005./717.
HU_{dyn}
\]
avec $\gamma =  1005/717 = 1.40167364$


%%%%%AUTEUR%%%%
\itemstyle{Auteur~:}
Le programme originel �tait �crit en Pascal, et nous a �t�
transmis par M. Helten (Research centre Juelich GmbH). La version
en C, aussi fid�le que possible � l'original, a �t� �crite par B. Piguet.

%%%%%REFERENCE%%%%
\itemstyle{R\'ef\'erences~:}
La tension de vapeur saturante est calcul�e avec
les formules 3.17 et 3.18 donn�es dans Sonntag (1982)
% : Formula 3.17 & 3.18
%  with regards to Liquid Water and Ice respectivily.

